\SixthArticle{11}
If we wished to exhaust all relations between three functions $F(\alpha,\beta,\gamma)$, $F(\alpha+\lambda,\beta+\mu,\gamma+\nu)$, $F(\alpha+\lambda',\beta+\mu',\gamma+\nu')$, in which $\lambda,\mu,\nu,\lambda',\mu',\nu'$ are either $=0$, $=+1$ or $=-1$, the number of formulas would rise to 325. Such a collection, at least of the simpler among these formulas, would not be useless. Here, however, it may suffice to have included just a few, which anyone will be able to prove for themselves without difficulty, either from the formulas of Article 7 or, if they prefer, in the same way as the first two of those were obtained in Article 8.
\[
\begin{aligned}&[16]\quad F(\alpha,\beta,\gamma)-F(\alpha,\beta,\gamma-1)\\&{}=-\frac{\alpha\beta x}{\gamma(\gamma-1)}F(\alpha+1,\beta+1,\gamma+1)
\end{aligned}
\]
\[
\begin{aligned}&[17]\quad F(\alpha,\beta+1,\gamma)-F(\alpha,\beta,\gamma)\\&{}=\frac{\alpha x}{\gamma}F(\alpha+1,\beta+1,\gamma+1)
\end{aligned}
\]
\[
\begin{aligned}&[18]\quad F(\alpha+1,\beta,\gamma)-F(\alpha,\beta,\gamma)\\&{}=\frac{\beta x}{\gamma}F(\alpha+1,\beta+1,\gamma+1)
\end{aligned}
\]
\[
\begin{aligned}&[19]\quad F(\alpha,\beta+1,\gamma+1)-F(\alpha,\beta,\gamma)\\&{}=\frac{\alpha(\gamma-\beta)x}{\gamma(\gamma+1)}F(\alpha+1,\beta+1,\gamma+2)
\end{aligned}
\]
\[
\begin{aligned}&[20]\quad F(\alpha+1,\beta,\gamma+1)-F(\alpha,\beta,\gamma)\\&{}=\frac{\beta(\gamma-\alpha)x}{\gamma(\gamma+1)}F(\alpha+1,\beta+1,\gamma+2)
\end{aligned}
\]
\[
\begin{aligned}&[21]\quad F(\alpha-1,\beta+1,\gamma)-F(\alpha,\beta,\gamma)\\&{}=\frac{(\alpha-\beta-1)x}{\gamma}F(\alpha,\beta+1,\gamma+1)
\end{aligned}
\]
\[
\begin{aligned}&[22]\quad F(\alpha+1,\beta-1,\gamma)-F(\alpha,\beta,\gamma)\\&{}=\frac{(\beta-\alpha-1)x}{\gamma}F(\alpha+1,\beta,\gamma+1)
\end{aligned}
\]
\[
\begin{aligned}&[23]\quad F(\alpha-1,\beta+1,\gamma)-F(\alpha+1,\beta-1,\gamma)\\&{}=\frac{(\alpha-\beta)x}{\gamma}F(\alpha+1,\beta+1,\gamma+1)
\end{aligned}
\]
\begin{center}\rule{.24\linewidth}{.3pt}\end{center}
\SourceLeaf{143}
\par\Needspace{9\baselineskip}\pdfbookmark[1]{Continued fractions}{sixth-section-2}\begin{center}\textbf{\textit{SECOND SECTION.}}\\[6pt]\textbf{\textit{Continued fractions.}}\end{center}
\SixthArticle{12}
Denote
\[
\frac{F(\alpha,\beta+1,\gamma+1,x)}{F(\alpha,\beta,\gamma,x)}\quad\hbox{by}\quad G(\alpha,\beta,\gamma,x)
\]
Then
\[
\frac{F(\alpha+1,\beta,\gamma+1,x)}{F(\alpha,\beta,\gamma,x)}=\frac{F(\beta,\alpha+1,\gamma+1,x)}{F(\beta,\alpha,\gamma,x)}=G(\beta,\alpha,\gamma,x)
\]
and hence dividing equation 19 by $F(\alpha,\beta+1,\gamma+1,x)$ gives
\[
1-\frac{1}{G(\alpha,\beta,\gamma,x)}=\frac{\alpha(\gamma-\beta)}{\gamma(\gamma+1)}x\,G(\beta+1,\alpha,\gamma+1,x)
\]
or
\[
[24]\qquad G(\alpha,\beta,\gamma,x)=\cfrac{1}{1-\frac{\alpha(\gamma-\beta)}{\gamma(\gamma+1)}x\,G(\beta+1,\alpha,\gamma+1,x)}
\]
and since likewise
\[
G(\beta+1,\alpha,\gamma+1,x)=\cfrac{1}{1-\frac{(\beta+1)(\gamma+1-\alpha)}{(\gamma+1)(\gamma+2)}x\,G(\alpha+1,\beta+1,\gamma+2,x)}
\]
etc., the following continued fraction results for $G(\alpha,\beta,\gamma,x)$:
\[
[25]\qquad\frac{F(\alpha,\beta+1,\gamma+1,x)}{F(\alpha,\beta,\gamma,x)}=\cfrac{1}{1-\cfrac{ax}{1-\cfrac{bx}{1-\cfrac{cx}{1-\cfrac{dx}{1-\hbox{etc.}}}}}}
\]
where
\[
\begin{array}{ll}
a=\frac{\alpha(\gamma-\beta)}{\gamma(\gamma+1)}&b=\frac{(\beta+1)(\gamma+1-\alpha)}{(\gamma+1)(\gamma+2)}\\[6pt]
c=\frac{(\alpha+1)(\gamma+1-\beta)}{(\gamma+2)(\gamma+3)}&d=\frac{(\beta+2)(\gamma+2-\alpha)}{(\gamma+3)(\gamma+4)}\\[6pt]
e=\frac{(\alpha+2)(\gamma+2-\beta)}{(\gamma+4)(\gamma+5)}&f=\frac{(\beta+3)(\gamma+3-\alpha)}{(\gamma+5)(\gamma+6)}
\end{array}
\]
etc.; its law of progression is evident.

Further, equations 17, 18, 21, 22 give
\SourceLeaf{144}
\[
\begin{array}{rl}
[26]&\frac{F(\alpha,\beta+1,\gamma,x)}{F(\alpha,\beta,\gamma,x)}=\cfrac{1}{1-\frac{\alpha x}{\gamma}G(\beta+1,\alpha,\gamma,x)}\\[12pt]
[27]&\frac{F(\alpha+1,\beta,\gamma,x)}{F(\alpha,\beta,\gamma,x)}=\cfrac{1}{1-\frac{\beta x}{\gamma}G(\alpha+1,\beta,\gamma,x)}\\[12pt]
[28]&\frac{F(\alpha-1,\beta+1,\gamma,x)}{F(\alpha,\beta,\gamma,x)}=\cfrac{1}{1-\frac{(\alpha-\beta-1)x}{\gamma}G(\beta+1,\alpha-1,\gamma,x)}\\[12pt]
[29]&\frac{F(\alpha+1,\beta-1,\gamma,x)}{F(\alpha,\beta,\gamma,x)}=\cfrac{1}{1-\frac{(\beta-\alpha-1)x}{\gamma}G(\alpha+1,\beta-1,\gamma,x)}
\end{array}
\]
Thus, substituting for the function $G$ its values as continued fractions gives as many new continued fractions.

Moreover, it is evident that the continued fraction in formula 25 terminates if any of the numbers $\alpha,\beta,\gamma-\alpha,\gamma-\beta$ is a negative integer, but otherwise continues indefinitely.
\SixthArticle{13}
The continued fractions obtained in the preceding article are of the greatest importance, and it can be asserted that hardly any continued fractions following an evident law have so far been found by analysts which are not included in ours as special cases. Particularly noteworthy is the case in which $\beta=0$ is set in formula 25. This gives $F(\alpha,\beta,\gamma,x)=1$, and so, on writing $\gamma-1$ for $\gamma$,
\[
\begin{aligned}
{}[30]\quad F(\alpha,1,\gamma)&=1+\frac{\alpha}{\gamma}x+\frac{\alpha(\alpha+1)}{\gamma(\gamma+1)}xx\\&{}+\frac{\alpha(\alpha+1)(\alpha+2)}{\gamma(\gamma+1)(\gamma+2)}x^3+\hbox{etc.}\\[9pt]
&=\cfrac{1}{1-\cfrac{ax}{1-\cfrac{bx}{1-\cfrac{cx}{1-\cfrac{dx}{1-\hbox{etc.}}}}}}
\end{aligned}
\]
where
\[
\begin{array}{ll}
a=\frac{\alpha}{\gamma}&b=\frac{\gamma-\alpha}{\gamma(\gamma+1)}\\[6pt]
c=\frac{(\alpha+1)\gamma}{(\gamma+1)(\gamma+2)}&d=\frac{2(\gamma+1-\alpha)}{(\gamma+2)(\gamma+3)}\\[6pt]
e=\frac{(\alpha+2)(\gamma+1)}{(\gamma+3)(\gamma+4)}&f=\frac{3(\gamma+2-\alpha)}{(\gamma+4)(\gamma+5)}
\end{array}
\]
\begin{center}etc.\end{center}
